% GENERATED FILE. Edit canonical lessons, then run npm run generate:books.
% canonical-source-hash: fnv1a64:eb2dc11755399af0
% canonical-lessons: SA-C202-pathikah, SA-C202-disa, SA-C202-purva, SA-C202-pascima, SA-C202-uttara, SA-R202-first-pass-words-from-chapters-197-199, SA-R202-second-pass-words-from-chapters-200-202

\chapter{Wayfarer, Direction, East, West, North}
\label{ch:sa-wayfarer-direction-east-west-north}
\hlchaptermodality{202}

\begin{chapteropening}
\textbf{By the end of this chapter:} \emph{I can say a wayfarer, a direction, the east, the west and the north.}

Complete the last lesson of chapter 202: I can say a wayfarer, a direction, the east, the west and the north.
\end{chapteropening}

\section[\texorpdfstring{\sk{पथिकः} (pathikaḥ)}{पथिकः (pathikaḥ)}]{\sk{पथिकः} (pathikaḥ) — a wayfarer}

\label{lesson:SA-C202-pathikah}



\begin{quote}
Before the new one: say the Sanskrit for a problem, then the Sanskrit for a solution.
\end{quote}

\subsection*{You'll want to know: \sk{पथिकः}}
\textbf{\sk{पथिकः}} — \emph{pathikaḥ} — \textquotedblleft{}a wayfarer\textquotedblright{}.

\begin{grammarlens}[title={What you've built}]
One more word for this chapter.
\end{grammarlens}

\subsection*{Your turn}
\begin{itemize}
\raggedright
  \item \emph{Say it:} \emph{pathikaḥ}
  \item \emph{Say it:} \emph{pathikaḥ}, once more
  \item \emph{Say it:} say \emph{pathikaḥ}
  \item \emph{Recall:} say the Sanskrit for a problem, then the Sanskrit for a solution, then say \emph{pathikaḥ} again
\end{itemize}

\begin{tcolorbox}[breakable,colback=teal!4,colframe=teal!35!black,title={Before you move on}]
What does \sk{पथिकः} mean? (\textquotedblleft{}A wayfarer\textquotedblright{}.) Say it once more.
\end{tcolorbox}
\section[\texorpdfstring{\sk{दिशा} (diśā)}{दिशा (diśā)}]{\sk{दिशा} (diśā) — a direction}

\label{lesson:SA-C202-disa}



\begin{quote}
Before the new one: say the Sanskrit for a solution, then the Sanskrit for a wayfarer.
\end{quote}

\subsection*{You'll want to know: \sk{दिशा}}
\textbf{\sk{दिशा}} — \emph{diśā} — \textquotedblleft{}a direction\textquotedblright{}.

\begin{grammarlens}[title={What you've built}]
One more word for this chapter.
\end{grammarlens}

\subsection*{Your turn}
\begin{itemize}
\raggedright
  \item \emph{Say it:} \emph{diśā}
  \item \emph{Say it:} \emph{diśā}, once more
  \item \emph{Say it:} say \emph{diśā}
  \item \emph{Recall:} say the Sanskrit for a solution, then the Sanskrit for a wayfarer, then say \emph{diśā} again
\end{itemize}

\begin{tcolorbox}[breakable,colback=teal!4,colframe=teal!35!black,title={Before you move on}]
What does \sk{दिशा} mean? (\textquotedblleft{}A direction\textquotedblright{}.) Say it once more.
\end{tcolorbox}
\section[\texorpdfstring{\sk{पूर्वा} (pūrvā)}{पूर्वा (pūrvā)}]{\sk{पूर्वा} (pūrvā) — the east}

\label{lesson:SA-C202-purva}



\begin{quote}
Before the new one: say the Sanskrit for a wayfarer, then the Sanskrit for a direction.
\end{quote}

\subsection*{You'll want to know: \sk{पूर्वा}}
\textbf{\sk{पूर्वा}} — \emph{pūrvā} — \textquotedblleft{}the east\textquotedblright{}.

\begin{grammarlens}[title={What you've built}]
One more word for this chapter.
\end{grammarlens}

\subsection*{Your turn}
\begin{itemize}
\raggedright
  \item \emph{Say it:} \emph{pūrvā}
  \item \emph{Say it:} \emph{pūrvā}, once more
  \item \emph{Say it:} say \emph{pūrvā}
  \item \emph{Recall:} say the Sanskrit for a wayfarer, then the Sanskrit for a direction, then say \emph{pūrvā} again
\end{itemize}

\begin{tcolorbox}[breakable,colback=teal!4,colframe=teal!35!black,title={Before you move on}]
What does \sk{पूर्वा} mean? (\textquotedblleft{}The east\textquotedblright{}.) Say it once more.
\end{tcolorbox}
\section[\texorpdfstring{\sk{पश्चिमा} (paścimā)}{पश्चिमा (paścimā)}]{\sk{पश्चिमा} (paścimā) — the west}

\label{lesson:SA-C202-pascima}



\begin{quote}
Before the new one: say the Sanskrit for a direction, then the Sanskrit for the east.
\end{quote}

\subsection*{You'll want to know: \sk{पश्चिमा}}
\textbf{\sk{पश्चिमा}} — \emph{paścimā} — \textquotedblleft{}the west\textquotedblright{}.

\begin{grammarlens}[title={What you've built}]
One more word for this chapter.
\end{grammarlens}

\subsection*{Your turn}
\begin{itemize}
\raggedright
  \item \emph{Say it:} \emph{paścimā}
  \item \emph{Say it:} \emph{paścimā}, once more
  \item \emph{Say it:} say \emph{paścimā}
  \item \emph{Recall:} say the Sanskrit for a direction, then the Sanskrit for the east, then say \emph{paścimā} again
\end{itemize}

\begin{tcolorbox}[breakable,colback=teal!4,colframe=teal!35!black,title={Before you move on}]
What does \sk{पश्चिमा} mean? (\textquotedblleft{}The west\textquotedblright{}.) Say it once more.
\end{tcolorbox}
\section[\texorpdfstring{\sk{उत्तरा} (uttarā)}{उत्तरा (uttarā)}]{\sk{उत्तरा} (uttarā) — the north}

\label{lesson:SA-C202-uttara}



\begin{quote}
Before the new one: say the Sanskrit for the east, then the Sanskrit for the west.
\end{quote}

\subsection*{You'll want to know: \sk{उत्तरा}}
\textbf{\sk{उत्तरा}} — \emph{uttarā} — \textquotedblleft{}the north\textquotedblright{}.

\begin{grammarlens}[title={What you've built}]
One more word for this chapter.
\end{grammarlens}

\subsection*{Your turn}
\begin{itemize}
\raggedright
  \item \emph{Say it:} \emph{uttarā}
  \item \emph{Say it:} \emph{uttarā}, once more
  \item \emph{Say it:} say \emph{uttarā}
  \item \emph{Recall:} say the Sanskrit for the east, then the Sanskrit for the west, then say \emph{uttarā} again
\end{itemize}

\begin{tcolorbox}[breakable,colback=teal!4,colframe=teal!35!black,title={Before you move on}]
What does \sk{उत्तरा} mean? (\textquotedblleft{}The north\textquotedblright{}.) Say it once more.
\end{tcolorbox}
\section[(dialogue)]{Words from chapters 197-199 — a first pass}

\label{lesson:SA-R202-first-pass-words-from-chapters-197-199}



\begin{quote}
Say the words for the west and the north.
\end{quote}

\subsection*{Your turn}
Say these aloud:

\begin{itemize}
\raggedright
  \item a bamboo grove, a rainbow, hail, dew, moonlight
  \item a deer, a leopard, a rhinoceros, a nest, an egg
  \item a sentence, a language, grammar, mathematics, science
\end{itemize}

\begin{tcolorbox}[breakable,colback=teal!4,colframe=teal!35!black,title={Before you move on}]
A bamboo grove? (\textbf{\sk{वेणुवनम्}}, \emph{veṇuvanam}.) A rainbow? (\textbf{\sk{इन्द्रधनुः}}, \emph{indradhanuḥ}.) Hail? (\textbf{\sk{करका}}, \emph{karakā}.) Dew? (\textbf{\sk{अवश्यायः}}, \emph{avaśyāyaḥ}.) Moonlight? (\textbf{\sk{ज्योत्स्ना}}, \emph{jyotsnā}.) A deer? (\textbf{\sk{हरिणः}}, \emph{hariṇaḥ}.) A leopard? (\textbf{\sk{चित्रकः}}, \emph{citrakaḥ}.) A rhinoceros? (\textbf{\sk{गण्डकः}}, \emph{gaṇḍakaḥ}.) A nest? (\textbf{\sk{नीडम्}}, \emph{nīḍam}.) An egg? (\textbf{\sk{अण्डम्}}, \emph{aṇḍam}.) A sentence? (\textbf{\sk{वाक्यम्}}, \emph{vākyam}.) A language? (\textbf{\sk{भाषा}}, \emph{bhāṣā}.) Grammar? (\textbf{\sk{व्याकरणम्}}, \emph{vyākaraṇam}.) Mathematics? (\textbf{\sk{गणितम्}}, \emph{gaṇitam}.) Science? (\textbf{\sk{विज्ञानम्}}, \emph{vijñānam}.)
\end{tcolorbox}
\section[(dialogue)]{Words from chapters 200-202 — a second pass}

\label{lesson:SA-R202-second-pass-words-from-chapters-200-202}



\begin{quote}
Say the words for reading and listening.
\end{quote}

\subsection*{Your turn}
Say these aloud:

\begin{itemize}
\raggedright
  \item reading, listening, a speech, a discussion, a dispute
  \item a plan, a goal, a way, a problem, a solution
  \item a wayfarer, a direction, the east, the west, the north
\end{itemize}

\begin{tcolorbox}[breakable,colback=teal!4,colframe=teal!35!black,title={Before you move on}]
Reading? (\textbf{\sk{पठनम्}}, \emph{paṭhanam}.) Listening? (\textbf{\sk{श्रवणम्}}, \emph{śravaṇam}.) A speech? (\textbf{\sk{भाषणम्}}, \emph{bhāṣaṇam}.) A discussion? (\textbf{\sk{चर्चा}}, \emph{carcā}.) A dispute? (\textbf{\sk{विवादः}}, \emph{vivādaḥ}.) A plan? (\textbf{\sk{योजना}}, \emph{yojanā}.) A goal? (\textbf{\sk{लक्ष्यम्}}, \emph{lakṣyam}.) A way? (\textbf{\sk{उपायः}}, \emph{upāyaḥ}.) A problem? (\textbf{\sk{समस्या}}, \emph{samasyā}.) A solution? (\textbf{\sk{समाधानम्}}, \emph{samādhānam}.) A wayfarer? (\textbf{\sk{पथिकः}}, \emph{pathikaḥ}.) A direction? (\textbf{\sk{दिशा}}, \emph{diśā}.) The east? (\textbf{\sk{पूर्वा}}, \emph{pūrvā}.) The west? (\textbf{\sk{पश्चिमा}}, \emph{paścimā}.) The north? (\textbf{\sk{उत्तरा}}, \emph{uttarā}.)
\end{tcolorbox}
